% Compilação recomendada: xelatex logica.tex (duas vezes).
\documentclass[10pt,aspectratio=169]{beamer}

% ------------------------------------------------------------
% Tema: Metropolis (limpo e elegante)
% ------------------------------------------------------------
\usetheme{metropolis}
\metroset{block=fill, sectionpage=none, subsectionpage=none, numbering=fraction}

\usepackage{iftex}
\ifPDFTeX\usepackage[utf8]{inputenc}\usepackage[T1]{fontenc}\else\usepackage{fontspec}\setsansfont{Latin Modern Sans}\fi
\usepackage{amsmath,amssymb}
\usepackage{mathtools}
\usepackage{tikz}
\usetikzlibrary{positioning,arrows.meta}
\usepackage{forest}

\usepackage{booktabs}
\usepackage{array}

\setbeamertemplate{frame footer}{Lógica para Computação}

% Paleta discreta
\definecolor{mDarkBlue}{HTML}{2E3A46}
\definecolor{mAccent}{HTML}{8C6D46}
\setbeamercolor{palette primary}{bg=mDarkBlue,fg=white}
\setbeamercolor{progress bar}{fg=mAccent,bg=mDarkBlue!20}
\setbeamercolor{alerted text}{fg=mAccent}

% Macro para referência ao livro-base, exibida discretamente no rodapé do título de seção
\newcommand{\refbook}[1]{{\footnotesize\color{mAccent}\textit{Ref.: #1}}}

\newcommand{\Prop}{\mathrm{Prop}}
\newcommand{\dest}[1]{\alert{#1}}

\title{Lógica para a Computação}
\subtitle{Resumo de aula: lógica proposicional clássica}
\author{}
\date{}

\begin{document}

\begin{frame}[plain]
\vfill
{\LARGE\bfseries Lógica para a Computação\par}
\medskip
{\large Resumo de aula: lógica proposicional clássica\par}
\bigskip
{\color{mAccent} Sintaxe, semântica e sistemas de prova\par}
\vfill
\end{frame}

\begin{frame}{Sumário}
  \small
  \begin{columns}[T]
    \begin{column}{0.5\textwidth}
      \begin{enumerate}
        \item Expressividade e sintaxe
        \item Dedução natural
        \item Semântica
        \item Modelos, satisfazibilidade e validade
        \item Tabela verdade e falsificação
        \item Consequência lógica
      \end{enumerate}
    \end{column}
    \begin{column}{0.5\textwidth}
      \begin{enumerate}
        \setcounter{enumi}{6}
        \item Tableaux semânticos
        \item Corretude e completude
        \item Equivalências
        \item Formas normais
        \item Resolução
      \end{enumerate}
    \end{column}
  \end{columns}
\end{frame}

% ==============================================================
\section{Expressividade e sintaxe}
% ==============================================================
\begin{frame}{Alfabeto da lógica proposicional}
  \begin{block}{Símbolos primitivos}
    \begin{itemize}
      \item Variáveis proposicionais: $p, q, r, \dots \in \mathrm{Prop}$
      \item Conectivos: $\neg,\ \land,\ \lor,\ \rightarrow,\ \leftrightarrow$
      \item Símbolos auxiliares: $(\,,\,)$
    \end{itemize}
  \end{block}
  \begin{alertblock}{Definição indutiva de fórmula (FBF)}
    \begin{enumerate}
      \item Toda variável $p\in\mathrm{Prop}$ é uma fórmula;
      \item Se $\varphi$ é fórmula, $\neg\varphi$ também é;
      \item Se $\varphi$ e $\psi$ são fórmulas, então $(\varphi\land\psi)$, $(\varphi\lor\psi)$, $(\varphi\rightarrow\psi)$ e $(\varphi\leftrightarrow\psi)$ também são;
      \item Nada mais é fórmula.
    \end{enumerate}
  \end{alertblock}
\end{frame}

\begin{frame}{Precedência e árvore sintática}
  \begin{columns}[T]
    \begin{column}{0.48\textwidth}
      \textbf{Convenção de precedência} (do mais forte ao mais fraco):
      \[ \neg \;>\; \land \;>\; \lor \;>\; \rightarrow \;>\; \leftrightarrow \]
      Isso permite escrever $p\lor q \rightarrow \neg r$ em vez de\\
      $((p\lor q)\rightarrow(\neg r))$.
    \end{column}
    \begin{column}{0.48\textwidth}
      \centering
      \textbf{Árvore de } $(p\lor q)\rightarrow \neg r$\par\medskip
      \begin{tikzpicture}[level distance=10mm, sibling distance=16mm,
          every node/.style={font=\small}]
        \node {$\rightarrow$}
          child { node {$\lor$}
            child { node {$p$} }
            child { node {$q$} } }
          child { node {$\neg$}
            child { node {$r$} } };
      \end{tikzpicture}
    \end{column}
  \end{columns}
\end{frame}

% ==============================================================

\begin{frame}{Linguagem, fórmulas e metalinguagem}

\begin{itemize}
\item $p,q,r$: variáveis proposicionais. Representam afirmações tratadas como unidades.
\item $\varphi,\psi,\chi$: metavariáveis. Representam fórmulas arbitrárias.
\item $\Gamma,\Delta$: conjuntos de fórmulas, usualmente usados como premissas.
\item Um \dest{literal} é uma variável $p$ ou sua negação $\neg p$.
\item $\bot$ e $\top$ abreviam $p\land\neg p$ e $p\lor\neg p$, respectivamente.
\end{itemize}
\begin{block}{Escopo e convenções}
Trabalhamos na lógica proposicional \textbf{clássica}, com dois valores.
Os conectivos são verofuncionais. Quantificadores e estrutura interna de predicados exigem outra linguagem.
\end{block}
Usaremos parênteses para evitar ambiguidades entre conectivos repetidos.
\end{frame}

\section{Dedução natural}
% ==============================================================

\begin{frame}{Sistema de dedução natural}
  \begin{itemize}
    \item Sistema \dest{sintático}: prova-se $\varphi$ manipulando símbolos, sem apelar a valores-verdade.
    \item Cada conectivo tem regras de \dest{introdução} (I) e \dest{eliminação} (E).
    \item Hipóteses podem ser abertas e depois \dest{descartadas} (fechadas) dentro de uma subprova.
  \end{itemize}
  \vspace{4mm}
  \begin{block}{Duas regras centrais}
    \[
      \dfrac{\varphi \quad \varphi\rightarrow\psi}{\psi}\ (\rightarrow\!E,\ \textit{Modus Ponens})
      \qquad\qquad
      \cfrac{\begin{array}{c}[\varphi]\\ \vdots \\ \psi\end{array}}{\varphi\rightarrow\psi}\ (\rightarrow\!I)
    \]
  \end{block}
\end{frame}


\begin{frame}{Exemplo de prova}

\textbf{Objetivo:} $\vdash (p\land q)\to p$.
\begin{center}
\renewcommand{\arraystretch}{1.5}
\begin{tabular}{cll}
\toprule
Linha & Fórmula & Justificativa\\\midrule
1 & $[p\land q]^1$ & Hipótese temporária\\
2 & $p$ & $\land E$ aplicada à linha 1\\
3 & $(p\land q)\to p$ & $\to I$, linhas 1--2, descarga de $1$\\
\bottomrule
\end{tabular}\par
\end{center}
A conclusão não depende de hipóteses abertas: é um \dest{teorema}.
\begin{block}{O que a descarga faz}
A subprova estabelece que, sob a hipótese $p\land q$, obtemos $p$.
A regra $\to I$ transforma essa dependência em uma implicação.
\end{block}
\end{frame}


% ==============================================================

\begin{frame}{Derivabilidade e hipóteses abertas}

\[\Gamma\vdash_S\varphi\]
Significa que existe uma derivação finita de $\varphi$ no sistema $S$, cujas hipóteses não descarregadas pertencem a $\Gamma$.
\begin{itemize}
\item \dest{Teorema}: $\vdash_S\varphi$, uma derivação sem premissas abertas.
\item \dest{Consistência sintática}: $\Gamma\nvdash_S\bot$.
\item Uma hipótese temporária pode desaparecer do conjunto de dependências por uma regra de descarga.
\end{itemize}
\begin{block}{Exemplo}
$\{p\land q\}\vdash p$ usa uma premissa.\\
$\vdash(p\land q)\to p$ expressa um teorema.
\end{block}
\end{frame}

\begin{frame}{Dedução natural: conjunção e disjunção}

\[\frac{\varphi\quad\psi}{\varphi\land\psi}\;(\land I)
\qquad\frac{\varphi\land\psi}{\varphi}\;(\land E_1)
\qquad\frac{\varphi\land\psi}{\psi}\;(\land E_2)\]
\[\frac{\varphi}{\varphi\lor\psi}\;(\lor I_1)
\qquad\frac{\psi}{\varphi\lor\psi}\;(\lor I_2)\]
\[\frac{\varphi\lor\psi\qquad
\begin{array}{c}[\varphi]^i\\\vdots\\\chi\end{array}\qquad
\begin{array}{c}[\psi]^j\\\vdots\\\chi\end{array}}
{\chi}\;(\lor E,i,j)\]
Na eliminação de $\lor$, a mesma conclusão deve seguir de \dest{cada alternativa}. As hipóteses marcadas são descarregadas.
\end{frame}

\begin{frame}{Dedução natural: negação e lógica clássica}

\[\frac{\begin{array}{c}[\varphi]^i\\\vdots\\\bot\end{array}}{\neg\varphi}\;(\neg I,i)
\qquad\frac{\varphi\quad\neg\varphi}{\bot}\;(\neg E)
\qquad\frac{\bot}{\psi}\;(\bot E)\]
\begin{block}{Uma regra especificamente clássica}
\[\frac{\neg\neg\varphi}{\varphi}\;(\mathrm{DNE})\]
A eliminação da dupla negação permite recuperar a lógica clássica sobre as regras intuicionistas usuais.
\end{block}
O bicondicional pode ser abreviação de $(\varphi\to\psi)\land(\psi\to\varphi)$. A completude depende das regras efetivamente adotadas.
\end{frame}

\section{Semântica}
% ==============================================================
\begin{frame}{Valorações}
  \begin{block}{Definição}
    Uma \dest{valoração} é uma função $v:\mathrm{Prop}\to\{V,F\}$, estendida a todas as fórmulas por:
  \end{block}
  \begin{center}
    \begin{tabular}{c|c}
      \toprule
      $\varphi$ & $v(\varphi)$ \\
      \midrule
      $\neg\alpha$ & $V$ sse $v(\alpha)=F$\\
      $\alpha\land\beta$ & $V$ sse $v(\alpha)=V$ e $v(\beta)=V$\\
      $\alpha\lor\beta$ & $V$ sse $v(\alpha)=V$ ou $v(\beta)=V$\\
      $\alpha\rightarrow\beta$ & $F$ sse $v(\alpha)=V$ e $v(\beta)=F$\\
      $\alpha\leftrightarrow\beta$ & $V$ sse $v(\alpha)=v(\beta)$\\
      \bottomrule
    \end{tabular}\par
  \end{center}
\end{frame}

\begin{frame}{Avaliação de uma fórmula}
  Seja $v(p)=V,\ v(q)=F,\ v(r)=V$. Avalie $\varphi = (p\rightarrow q)\lor r$:
  \[
    v(p\rightarrow q) = F \qquad\Rightarrow\qquad v(\varphi) = F \lor v(r) = F\lor V = V
  \]
  \begin{alertblock}{Ideia-chave}
    A semântica atribui \dest{significado} (valor-verdade) às fórmulas sintáticas — independentemente de como elas foram provadas.
  \end{alertblock}
\end{frame}

% ==============================================================
\section{Modelos, satisfazibilidade e validade}
% ==============================================================

\begin{frame}{Classificação das fórmulas}

\small
\renewcommand{\arraystretch}{1.45}
\begin{tabular}{p{2.7cm}p{7.9cm}}
\toprule
Noção & Definição\\\midrule
Satisfação / modelo & $v\models\varphi$ significa $v(\varphi)=V$. A valoração $v$ é modelo de $\varphi$.\\
Satisfazibilidade & $\varphi$ é satisfazível se $\exists v\;(v\models\varphi)$.\\
Validade / tautologia & $\models\varphi$ significa $\forall v\;(v\models\varphi)$.\\
Insatisfazibilidade & $\varphi$ é insatisfazível se $\forall v\;(v(\varphi)=F)$.\\
Contingência & Existem $v,w$ com $v(\varphi)=V$ e $w(\varphi)=F$.\\
\bottomrule
\end{tabular}\par
\medskip
Toda fórmula válida é satisfazível. Uma fórmula satisfazível pode ser contingente.
\end{frame}


\begin{frame}{Exemplo}
  \[ \varphi = p \lor \neg p \quad(\text{válida}) \qquad\qquad \psi = p\land \neg p \quad(\text{insatisfazível}) \]
  \[ \chi = p \rightarrow q \quad(\text{satisfazível, mas \emph{não} válida}) \]
  \begin{itemize}
    \item $\varphi$ é verdadeira em \emph{toda} valoração $\Rightarrow$ tautologia.
    \item $\chi$ é falsa apenas quando $v(p)=V,\,v(q)=F$: tem modelos e "não modelos".
  \end{itemize}
\end{frame}

% ==============================================================

\begin{frame}{Satisfazibilidade de conjuntos de fórmulas}

\[v\models\Gamma\quad\Longleftrightarrow\quad
\forall\gamma\in\Gamma\;(v\models\gamma).\]
$\Gamma$ é \dest{satisfazível} se existe \textbf{uma mesma valoração} que satisfaz todas as suas fórmulas.
\begin{block}{Cuidado com a ordem dos quantificadores}
$p$ e $\neg p$ são individualmente satisfazíveis, mas $\{p,\neg p\}$ é insatisfazível.
\end{block}
\[\models\varphi\iff\{\neg\varphi\}\text{ é insatisfazível}.\]
\[\varphi\text{ é contingente}\iff
\varphi\text{ e }\neg\varphi\text{ são satisfazíveis}.\]
O conjunto vazio é satisfeito por toda valoração.
\end{frame}

\section{Tabela verdade e falsificação}
% ==============================================================
\begin{frame}{Tabela verdade: método de decisão}
  Para $n$ variáveis, há $2^n$ valorações possíveis — construir todas verifica validade/satisfazibilidade por força bruta.
  \begin{center}
    \small
    \begin{tabular}{cc|c}
      \toprule
      $p$ & $q$ & $p\rightarrow q$\\
      \midrule
      V & V & V \\
      V & F & F \\
      F & V & V \\
      F & F & V \\
      \bottomrule
    \end{tabular}\par
  \end{center}
  \dest{Limitação}: o número de linhas cresce exponencialmente com o número de variáveis.
\end{frame}

\begin{frame}{Método da falsificação}
  \begin{block}{Ideia}
    Para testar se $\models\varphi$, tenta-se \dest{construir} uma valoração que torne $\varphi$ falsa. Se toda tentativa gera contradição, $\varphi$ é válida.
  \end{block}
  \textbf{Exemplo:} mostrar que $\models (p\rightarrow q)\lor(q\rightarrow p)$.
  \begin{itemize}
    \item Suponha $v\big((p\to q)\lor(q\to p)\big)=F$;
    \item Então $v(p\to q)=F$ \emph{e} $v(q\to p)=F$;
    \item Da primeira: $v(p)=V,v(q)=F$; da segunda: $v(q)=V,v(p)=F$ — \dest{contradição};
    \item Logo não há falsificação possível: a fórmula é válida.
  \end{itemize}
\end{frame}

% ==============================================================

\begin{frame}{Tabela-verdade do exemplo central}

\begin{center}
\renewcommand{\arraystretch}{1.45}
\begin{tabular}{cc|c|c|c}
\toprule
$p$&$q$&$p\land q$&$(p\land q)\to p$&$\neg((p\land q)\to p)$\\\midrule
V&V&V&V&F\\
V&F&F&V&F\\
F&V&F&V&F\\
F&F&F&V&F\\\bottomrule
\end{tabular}\par
\end{center}
A implicação é verdadeira em todas as linhas. Sua negação é insatisfazível.
\begin{block}{Ligação com o tableau}
O tableau tenta construir uma linha com a última coluna verdadeira.
O fechamento demonstra que essa linha não existe.
\end{block}
\end{frame}

\section{Consequência lógica}
% ==============================================================
\begin{frame}{Consequência lógica}
  \begin{block}{Definição}
    $\Gamma \models \varphi$ (``$\varphi$ é consequência lógica de $\Gamma$'') se \dest{toda} valoração que satisfaz todas as fórmulas de $\Gamma$ também satisfaz $\varphi$.
  \end{block}
  \begin{itemize}
    \item Caso $\Gamma=\emptyset$: recai na validade, $\models\varphi$.
    \item Para $\Gamma=\{\gamma_1,\dots,\gamma_n\}$ finito:
    \[
      \Gamma\models\varphi \quad\Longleftrightarrow\quad \models (\gamma_1\land\cdots\land\gamma_n)\rightarrow\varphi
    \]
  \end{itemize}
  \textbf{Exemplo:} $\{p\rightarrow q,\, p\}\models q$ (Modus Ponens é válido semanticamente).
\end{frame}

% ==============================================================

\begin{frame}{Contramodelos e redução à insatisfazibilidade}

\[\Gamma\not\models\varphi\iff
\exists v\;(v\models\Gamma\ \text{e}\ v(\varphi)=F).\]
Essa valoração é um \dest{contramodelo do argumento}.
\[\boxed{\Gamma\models\varphi\iff
\Gamma\cup\{\neg\varphi\}\text{ é insatisfazível}}\]
\textbf{Exemplo inválido:} $\{p\to q,q\}\not\models p$.\\
O contramodelo $v(p)=F$, $v(q)=V$ satisfaz ambas as premissas e falsifica a conclusão.
\begin{block}{Premissas insatisfazíveis}
Se $\Gamma$ não tem modelos, então $\Gamma\models\varphi$ para toda $\varphi$.
Não existe contramodelo com todas as premissas verdadeiras.
\end{block}
\end{frame}

\section{Tableaux semânticos}
% ==============================================================

\begin{frame}{Tableaux: ideia geral}

Um tableau organiza a busca por uma valoração que satisfaça as fórmulas da raiz.
\begin{itemize}
\item Fórmulas no \dest{mesmo ramo} devem valer simultaneamente.
\item Uma \dest{ramificação} representa alternativas semânticas.
\item Um ramo fecha ($\times$) ao conter $\psi$ e $\neg\psi$. Quando $\psi$ é atômica, temos literais complementares.
\item O tableau fecha quando \dest{todos} os seus ramos fecham.
\end{itemize}
\begin{block}{Três tarefas, o mesmo método}
Satisfazibilidade de $\Gamma$: raiz $\Gamma$.\\
Validade de $\varphi$: raiz $\neg\varphi$.\\
Consequência $\Gamma\models\varphi$: raiz $\Gamma\cup\{\neg\varphi\}$.
\end{block}
\end{frame}




\begin{frame}{Tableaux: regras sem ramificação}

\small
\renewcommand{\arraystretch}{1.65}
\begin{center}
\begin{tabular}{ll}
\toprule
Fórmula no ramo & Acrescentar ao mesmo ramo\\\midrule
$\neg\neg\varphi$ & $\varphi$\\
$\varphi\land\psi$ & $\varphi$ e $\psi$\\
$\neg(\varphi\lor\psi)$ & $\neg\varphi$ e $\neg\psi$\\
$\neg(\varphi\to\psi)$ & $\varphi$ e $\neg\psi$\\\bottomrule
\end{tabular}\par
\end{center}
\begin{block}{Por que a implicação negada não ramifica?}
$\varphi\to\psi$ é falsa exatamente quando $\varphi$ é verdadeira \textbf{e} $\psi$ é falsa. As duas exigências são simultâneas.
\end{block}
\end{frame}

\begin{frame}{Tableaux: regras com ramificação}

\small
\renewcommand{\arraystretch}{1.6}
\begin{tabular}{lll}
\toprule
Fórmula & Ramo esquerdo & Ramo direito\\\midrule
$\varphi\lor\psi$ & $\varphi$ & $\psi$\\
$\neg(\varphi\land\psi)$ & $\neg\varphi$ & $\neg\psi$\\
$\varphi\to\psi$ & $\neg\varphi$ & $\psi$\\
$\varphi\leftrightarrow\psi$ & $\varphi,\psi$ & $\neg\varphi,\neg\psi$\\
$\neg(\varphi\leftrightarrow\psi)$ & $\varphi,\neg\psi$ & $\neg\varphi,\psi$\\\bottomrule
\end{tabular}\par
\medskip
Cada ramo filho \dest{herda todas as fórmulas do caminho anterior}.\\
Uma vírgula em uma célula indica fórmulas no mesmo ramo.\\
Basta uma alternativa satisfazível para manter o conjunto satisfazível.
\end{frame}
\begin{frame}{Exemplo: tableau para $(p\land q)\rightarrow p$}

\begin{columns}[T]
\begin{column}{.48\textwidth}
\centering
\begin{forest}
for tree={grow'=south, edge={-}, l sep=2mm, inner sep=1.5pt, font=\small}
[{$\neg((p\land q)\to p)$}
 [{$p\land q$}
  [{$\neg p$}
   [{$p$}
    [{$q$}
     [{$\times\quad p,\neg p$}]]]]]]
\end{forest}
\end{column}
\begin{column}{.50\textwidth}
\small
\textbf{1. Negar a fórmula a testar.}\\
Procuramos um contramodelo da implicação.
\medskip

\textbf{2. Expandir a implicação negada.}\\
$\neg(A\to B)$ exige $A$ e $\neg B$ no mesmo ramo.
\medskip

\textbf{3. Expandir a conjunção.}\\
$p\land q$ exige $p$ e $q$ no mesmo ramo.
\medskip

\textbf{4. Fechar o único ramo.}\\
$p$ e $\neg p$ são incompatíveis.
\end{column}
\end{columns}
\medskip
\dest{Logo, $\models(p\land q)\to p$.} Não há ramificação neste exemplo.
\end{frame}


% ==============================================================

\begin{frame}{Exemplo ramificado: modus ponens}

\textbf{Objetivo:} $\{p\to q,p\}\models q$. Raiz: $p\to q,p,\neg q$.
\begin{center}
\begin{forest}
for tree={grow'=south, l sep=5mm, s sep=22mm, font=\normalsize}
[{$p\to q,\quad p,\quad\neg q$}
 [{$\neg p$} [{$\times\quad p,\neg p$}]]
 [{$q$} [{$\times\quad q,\neg q$}]]]
\end{forest}
\end{center}
A regra de $\to$ cria duas alternativas. Ambas entram em conflito com fórmulas herdadas da raiz.
\medskip

Todos os ramos fecham. Portanto, não há contramodelo do argumento.
\end{frame}

\begin{frame}{Ramo aberto, saturação e leitura de modelo}

Um ramo é \dest{saturado} quando todas as decomposições exigidas pelas regras já foram realizadas nele.
\begin{itemize}
\item Um ramo aberto ainda incompleto \textbf{não garante} satisfazibilidade.
\item Num ramo aberto saturado, atribua $V$ aos átomos positivos e $F$ aos negados. Átomos ausentes podem receber qualquer valor.
\end{itemize}
\begin{block}{Exemplo: $p\to q$ não é válida}
A raiz $\neg(p\to q)$ produz $p$ e $\neg q$ no mesmo ramo.\\
O ramo está aberto e saturado. O modelo da raiz é $v(p)=V$, $v(q)=F$, um contramodelo da fórmula original.
\end{block}
Para uma raiz proposicional finita, expandir cada ocorrência apenas uma vez por ramo produz uma árvore finita.
\end{frame}

\begin{frame}{Tableaux: roteiro de aplicação}

\begin{enumerate}
\item Escolher a raiz de acordo com a tarefa: satisfazibilidade, validade ou consequência.
\item Aplicar as regras sem ramificação antes das ramificações, quando conveniente.
\item Continuar as decomposições em cada ramo aberto, preservando as fórmulas herdadas.
\item Fechar ramos com fórmulas complementares.
\item Se todos fecham, a raiz é insatisfazível. Se sobra ramo aberto saturado, extrair um modelo da raiz.
\end{enumerate}
\begin{alertblock}{Erro frequente}
Fórmulas em ramos diferentes não fecham um ramo. A contradição precisa ocorrer no mesmo caminho da raiz até uma folha.
\end{alertblock}
\end{frame}

\section{Corretude e completude}
% ==============================================================

\begin{frame}{Corretude e completude}

Fixe um cálculo $S$ para a lógica proposicional clássica.
\begin{block}{Corretude (\textit{soundness})}
\[\Gamma\vdash_S\varphi\quad\Longrightarrow\quad\Gamma\models\varphi.\]
Toda conclusão derivável preserva a verdade das premissas.
\end{block}
\begin{block}{Completude semântica forte}
\[\Gamma\models\varphi\quad\Longrightarrow\quad\Gamma\vdash_S\varphi.\]
Toda consequência semântica admite uma derivação no cálculo.
\end{block}
Para $\Gamma=\varnothing$: $\vdash_S\varphi\iff\models\varphi$.
Isso relaciona \dest{demonstrabilidade e validade}, não verdade em uma valoração particular.
\end{frame}


% ==============================================================

\begin{frame}{Por que tableaux são corretos e completos?}

\small
\textbf{Corretude refutacional.} Se há um tableau fechado para $\Sigma$, então $\Sigma$ é insatisfazível.
\begin{itemize}
\item As regras preservam um modelo em pelo menos uma extensão do ramo.
\item Um ramo fechado não admite modelo. Logo, um modelo da raiz impediria o fechamento de todos os ramos.
\end{itemize}
\textbf{Completude refutacional.} Se $\Sigma$ é insatisfazível, existe um tableau fechado para $\Sigma$.
\begin{itemize}
\item Para $\Sigma$ finito, uma expansão exaustiva termina.
\item Um ramo aberto saturado forneceria um modelo, por indução na estrutura das fórmulas. A insatisfazibilidade impede tal ramo.
\end{itemize}
Para conjuntos infinitos, a passagem a uma refutação finita usa a compacidade da lógica proposicional.
\end{frame}

\begin{frame}{Completude, consistência e decidibilidade}

\small
\begin{itemize}
\item \dest{Completude fraca}: $\models\varphi\Rightarrow\vdash_S\varphi$.
\item \dest{Completude forte}: $\Gamma\models\varphi\Rightarrow\Gamma\vdash_S\varphi$, inclusive para $\Gamma$ infinito.
\item Num cálculo clássico correto e fortemente completo:
\[\Gamma\nvdash_S\bot\iff\Gamma\text{ é satisfazível}.\]
\item \dest{Decidibilidade}: existe um algoritmo que termina para toda entrada e responde corretamente ao problema considerado.
\end{itemize}
Validade e satisfazibilidade proposicionais de entradas finitas são decidíveis por tabelas-verdade ou tableaux exaustivos.
\begin{alertblock}{Distinção}
Completude de um cálculo não implica, por si só, decidibilidade. Também não significa que uma teoria decida toda fórmula por $\varphi$ ou $\neg\varphi$.
\end{alertblock}
\end{frame}

\section{Equivalências}
% ==============================================================
\begin{frame}{Equivalência lógica}
  \begin{block}{Definição}
    $\varphi \equiv \psi \iff\ \models \varphi\leftrightarrow\psi$ (mesmo valor-verdade em toda valoração).
  \end{block}
  \begin{columns}[T]
    \begin{column}{0.5\textwidth}
      \small
      \begin{align*}
        \neg\neg p &\equiv p \\
        p\land q &\equiv q\land p \\
        p\lor(q\lor r) &\equiv (p\lor q)\lor r \\
        p\to q &\equiv \neg p \lor q
      \end{align*}
    \end{column}
    \begin{column}{0.5\textwidth}
      \small
      \begin{align*}
        \neg(p\land q) &\equiv \neg p \lor \neg q \quad\text{(De Morgan)}\\
        \neg(p\lor q) &\equiv \neg p \land \neg q \quad\text{(De Morgan)}\\
        p\land(q\lor r) &\equiv (p\land q)\lor(p\land r)\\
        p\to q &\equiv \neg q \to \neg p \quad\text{(contrapositiva)}
      \end{align*}
    \end{column}
  \end{columns}
\end{frame}

% ==============================================================

\begin{frame}{Equivalência, implicação e consequência}

\renewcommand{\arraystretch}{1.55}
\begin{tabular}{ll}
\toprule
Expressão & Natureza\\\midrule
$\varphi\to\psi$ & Fórmula da linguagem\\
$\varphi\leftrightarrow\psi$ & Fórmula da linguagem\\
$\varphi\equiv\psi$ & Relação de equivalência semântica\\
$\Gamma\models\varphi$ & Relação de consequência semântica\\
$\Gamma\vdash_S\varphi$ & Relação de derivabilidade sintática\\\bottomrule
\end{tabular}\par
\[\varphi\equiv\psi\iff
(\varphi\models\psi\ \text{e}\ \psi\models\varphi).\]
A equivalência lógica permite substituir uma subfórmula por outra equivalente em qualquer fórmula proposicional.
\end{frame}

\section{Formas normais}
% ==============================================================
\begin{frame}{FNC e FND}
  \begin{block}{Formas normais}
    \begin{itemize}
      \item \dest{FND} (disjuntiva): $\bigvee_i \big(\bigwedge_j \ell_{ij}\big)$ — disjunção de conjunções de literais.
      \item \dest{FNC} (conjuntiva): $\bigwedge_i \big(\bigvee_j \ell_{ij}\big)$ — conjunção de disjunções de literais (cláusulas).
    \end{itemize}
  \end{block}
  \begin{alertblock}{Algoritmo de conversão para FNC}
    \begin{enumerate}
      \item Eliminar $\leftrightarrow$ e $\rightarrow$ (usando equivalências);
      \item Empurrar $\neg$ para dentro até literais (De Morgan, dupla negação);
      \item Distribuir $\lor$ sobre $\land$ até obter cláusulas.
    \end{enumerate}
  \end{alertblock}
\end{frame}

\begin{frame}{Exemplo de conversão}
  \[
    p \rightarrow (q\land r)
    \ \equiv\ \neg p \lor (q\land r)
    \ \equiv\ \underbrace{(\neg p\lor q)\land(\neg p\lor r)}_{\text{FNC}}
  \]
  \begin{itemize}
    \item Passo 1: elimina-se $\rightarrow$;
    \item Passo 2: negação já está aplicada a uma variável proposicional;
    \item Passo 3: distribui-se $\lor$ sobre $\land$, obtendo duas cláusulas.
  \end{itemize}
\end{frame}

% ==============================================================

\begin{frame}{Formas normais: exemplos e cuidados}

\begin{itemize}
\item \dest{FNN}: negações apenas diante de átomos, usando $\land$ e $\lor$.
\item \dest{FNC}: $(p\lor\neg q)\land(q\lor r)$.
\item \dest{FND}: $(p\land q)\lor(\neg p\land r)$.
\end{itemize}
Uma cláusula é uma disjunção de literais. Uma cláusula unitária contém um único literal.
\begin{block}{Equivalência e equissatisfazibilidade}
Conversões por equivalências preservam todos os modelos na mesma linguagem, mas podem aumentar exponencialmente a fórmula.\\
Transformações como a de Tseitin usam variáveis auxiliares e preservam satisfazibilidade. Não se deve confundir isso com equivalência na linguagem ampliada.
\end{block}
\end{frame}

\section{Resolução}
% ==============================================================
\begin{frame}{Regra de resolução}
  \begin{block}{Regra}
    A partir de duas cláusulas com um literal complementar:
    \[
      \dfrac{\ (A \lor p)\qquad (B\lor\neg p)\ }{A\lor B}\ (\text{resolvente})
    \]
  \end{block}
  \begin{itemize}
    \item Opera apenas sobre \dest{cláusulas} (fórmulas em FNC).
    \item \dest{Refutação por resolução}: para provar $\Gamma\models\varphi$, nega-se $\varphi$, converte-se $\Gamma\cup\{\neg\varphi\}$ para FNC e aplica-se resolução até derivar a \dest{cláusula vazia} $\Box$ (contradição).
  \end{itemize}
\end{frame}

\begin{frame}{Exemplo de refutação}
  Mostrar $\{p\rightarrow q,\ p\}\models q$. Cláusulas (negando $q$):
  \[
    \{\neg p \lor q\},\qquad \{p\},\qquad \{\neg q\}
  \]
  \begin{enumerate}
    \item Resolver $\{\neg p\lor q\}$ com $\{p\}$ $\Rightarrow$ $\{q\}$;
    \item Resolver $\{q\}$ com $\{\neg q\}$ $\Rightarrow$ $\Box$ (cláusula vazia).
  \end{enumerate}
  Contradição encontrada $\Rightarrow$ o conjunto original com $\neg q$ é insatisfazível $\Rightarrow$ $q$ é consequência lógica.
\end{frame}

% ==============================================================

\begin{frame}{Resolução: alcance e condição de término}

\textbf{Corretude:} as cláusulas parentais implicam o resolvente.
\medskip

\textbf{Completude refutacional:} um conjunto finito de cláusulas é insatisfazível se, e somente se, a resolução pode derivar $\Box$.
\begin{itemize}
\item $\Box$ é a cláusula vazia: disjunção sem literais, sempre falsa.
\item O conjunto vazio de cláusulas representa uma conjunção vazia, sempre verdadeira.
\item Não encontrar $\Box$ em algumas tentativas não basta para concluir satisfazibilidade.
\item É preciso saturar o conjunto: adicionar todos os resolventes novos até obter $\Box$ ou não haver novas cláusulas.
\end{itemize}
Há finitas cláusulas distintas sobre um conjunto finito de átomos. Com eliminação de duplicatas, essa saturação termina.
\end{frame}

\begin{frame}{Quadro de consulta}

\small
\renewcommand{\arraystretch}{1.6}
\begin{tabular}{p{3.9cm}p{6.8cm}}
\toprule
Pergunta & Critério\\\midrule
$\varphi$ é satisfazível? & Alguma valoração satisfaz $\varphi$.\\
$\varphi$ é válida? & Toda valoração satisfaz $\varphi$.\\
$\Gamma\models\varphi$? & $\Gamma\cup\{\neg\varphi\}$ é insatisfazível.\\
$\Gamma\not\models\varphi$? & Há modelo de $\Gamma$ que falsifica $\varphi$.\\
Tableau fechado? & Todas as alternativas da raiz são impossíveis.\\
Ramo aberto saturado? & Permite obter um modelo da raiz.\\
Cálculo correto e completo? & $\Gamma\vdash_S\varphi\iff\Gamma\models\varphi$.\\\bottomrule
\end{tabular}\par
\end{frame}

\begin{frame}{Exercícios de revisão}

\begin{enumerate}
\item Classifique $p\lor\neg p$, $p\land\neg p$ e $p\to q$.
\item Prove $\vdash(p\land q)\to q$ por dedução natural.
\item Construa um tableau para testar $(p\lor q)\to p$.
\item Teste $\{p\to q,\neg q\}\models\neg p$ por tableau.
\item Converta $\neg(p\to(q\land r))$ para FNC e para FND.
\item Refute por resolução $\{p\lor q,\neg p,\neg q\}$.
\item Explique por que um ramo aberto não saturado é inconclusivo.
\end{enumerate}
\end{frame}

\begin{frame}{Respostas e indicações}

\small
\begin{enumerate}
\item Tautologia, contradição e contingência, respectivamente.
\item Hipótese $p\land q$, seguida de $\land E_2$ e descarga com $\to I$.
\item Não é válida. A raiz produz $p\lor q,\neg p$. O ramo com $p$ fecha e o ramo com $q$ dá $p=F,q=V$.
\item A raiz $p\to q,\neg q,\neg\neg p$ produz $p$. Os ramos $\neg p$ e $q$ fecham. A consequência é válida.
\item FNC: $p\land(\neg q\lor\neg r)$. FND: $(p\land\neg q)\lor(p\land\neg r)$.
\item Resolver $p\lor q$ com $\neg p$ dá $q$. Resolver $q$ com $\neg q$ dá $\Box$.
\item Uma expansão pendente pode revelar uma contradição. Por exemplo, $p\land\neg p$ antes de expandir a conjunção.
\end{enumerate}
\end{frame}


\end{document}
